\begin{ExerciseBlock}
\ExerciseHeading{05}{Fallende Masse}{rigid_bodies_falling_mass_and_pulley}
In der Abbildung ist eine Masse $m = 500\,$g an ein Seil gebunden, das über eine Rolle mit Radius $R = 30,0\,$cm gewickelt ist. Über die Dichteverteilung innerhalb der Rolle ist Ihnen nichts bekannt. Die Masse fällt in $T=0,782\,$s genau $1$m.

\begin{enumerate}
	\item Berechnen Sie das Trägheitsmoment $J$ der Rolle.

	\item Berechnen Sie die Winkelgeschwindigkeit $\omega$ der anfänglich ruhenden Rolle, wenn die Masse um 1m gefallen ist.

\end{enumerate}

\begin{center}
	\CachedTikzFigure{content/resources/chapters/rigid_bodies/tikz/falling_mass_and_pulley.tex}
\end{center}

Hinweis: Nehmen Sie an, dass die Rolle einem Vollzylinder entspricht.
\end{ExerciseBlock}
\UseExerciseSolution{rigid_bodies_falling_mass_and_pulley}
